% !TeX program = lualatex
% Plantilla: RESUMEN formalizado (skill «apuntes-formalizados»)
\documentclass[10pt,a4paper]{article}

\newcommand{\autor}{Autor}
\newcommand{\titulocorto}{Tema $\cdot$ Resumen}

\usepackage[resumen]{apuntes-oscuro}
\hypersetup{pdftitle={Tema. Resumen}, pdfauthor={\autor}}

% Notación propia del tema
% \newcommand{\Lop}{\mathcal{L}}

\begin{document}
\thispagestyle{empty}
\titular{Título del tema}{Resumen}{desarrollos simbólicos}
{\small\color{muted}\autor\hfill Notación: $y'\equiv\dd y/\dd x$.}

\section{Primera sección}
Partimos de \dots{} y separamos \dots:
\begin{equation}
a=b\;\Longrightarrow\;\res{c=d}
\end{equation}

\begin{teorema}[Nombre]
\begin{equation}\label{eq:clave}
\res{f=g}
\end{equation}
\end{teorema}

\begin{demostracion}
\begin{enumerate}
\item \emph{Paso}: $a=b$.
\item Sustituyendo: $b\overset{\eqref{eq:clave}}{=}c$.
\end{enumerate}
\end{demostracion}

\end{document}
